%% =================================================================
%% Wei, Mei, Zhao, Xiao, Sun, Zhang, Guo.
%% "Nondestructively identifying the mechanical behavior of soft
%%  tissues using surface deformation with an explicit inverse
%%  approach."
%% Appl. Math. Modelling 134 (2024) 126--147.
%%
%% Reconstruction of page 11 (printed p. 136) as study notes.
%% This page introduces the first TABLE (Table 1) — a new content
%% variation relative to pages 1--10. It also contains Fig. 11 at
%% the top and the start of "Case 2: A layered ring structure"
%% at the bottom.
%% Prose: extracted verbatim from the PDF text layer via pdftotext.
%% Math:  none on this page.
%% Table 1: reproduced from the visual Read of the PDF page, using
%%          booktabs for the rule structure the original uses.
%% Figure 11: placeholder box with caption + visual description.
%% =================================================================

\documentclass[11pt]{article}

\usepackage[margin=1in]{geometry}
\usepackage{amsmath, amssymb}
\usepackage{mathrsfs}
\usepackage{bm}
\usepackage{xcolor}
\usepackage{fancyhdr}
\usepackage{booktabs}
\usepackage{multirow}
\usepackage{array}

\pagestyle{fancy}
\fancyhf{}
\lhead{\small Y.~Mei et al.}
\rhead{\small Applied Mathematical Modelling 134 (2024) 126--147}
\cfoot{\small 136 $\to$ \arabic{page}}
\renewcommand{\headrulewidth}{0.4pt}

\setcounter{page}{1}

\begin{document}

% ---------- Figure 11 placeholder ----------
\noindent
\begin{center}
\fbox{\parbox{0.92\textwidth}{\centering\vspace{1em}
\textbf{[Figure 11 --- placeholder]}\\[0.8em]
Six-panel side-by-side comparison of the \textbf{Explicit Inverse
Method} (left column) vs.\ \textbf{Implicit Inverse Method} (right
column) on \textbf{Sample 1} (homogeneous, target SM $= 3.00$\,MPa),
with \textbf{1\,\% noise} added to the measured surface
displacements. Three rows correspond to 3, 6, and 9 displacement
datasets.\\[0.4em]
\textbf{Explicit (left column)}:
(a) 3 fields: SM $\in [3.0013, 3.0013]$\,MPa, uniform red;
(b) 6 fields: SM $\in [3.0000, 3.0001]$\,MPa, visible green/yellow
patches emerge (non-monotone: noisier than (a) or (c));
(c) 9 fields: SM $\in [2.9997, 2.9997]$\,MPa, uniform red.\\[0.3em]
\textbf{Implicit (right column)}:
(d) 3 fields: SM $\in [2.9981, 3.0083]$\,MPa, visible blue patch
along the left edge;
(e) 6 fields: SM $\in [2.9978, 3.0006]$\,MPa, small central blue
blob;
(f) 9 fields: SM $\in [2.9995, 2.9998]$\,MPa, small blue patches
remain near the edges.\\[0.3em]
\textbf{Takeaway}: under 1\,\% noise, both methods stay
quantitatively close to the true 3.00\,MPa value ($\le 0.01$\,MPa
deviation). The explicit method's 6-field result is the noisiest
explicit panel on this page --- an artifact not discussed in the
surrounding prose.
\vspace{1em}}}
\end{center}
\vspace{0.3em}
\begin{center}
\textbf{Fig. 11.} The reconstructed results with varying 3, 6 and 9
surface displacement datasets corresponding to Sample 1 (1\,\% noise
is introduced into the measured data, Unit: MPa).
\end{center}

\vspace{1em}

% ---------- Table 1 ----------
\noindent\textbf{Table 1}\\[0.2em]
{\small Comparison of the average relative error in the shear
modulus distribution obtained by implicit inverse and explicit
inverse methods for the layered structure.}
\vspace{0.3em}

\begin{center}
\small
\begin{tabular}{@{}ll ccc ccc@{}}
\toprule
\multicolumn{2}{l}{Inverse method}
  & \multicolumn{3}{c}{Explicit inverse method}
  & \multicolumn{3}{c}{Implicit inverse method} \\
\cmidrule(lr){3-5}\cmidrule(lr){6-8}
\multicolumn{2}{l}{Number of surface displacement datasets}
  & 3 & 6 & 9 & 3 & 6 & 9 \\
\midrule
\multirow{3}{*}{1\,\% noise}
  & Sample 1 & 0.04\,\%  & 0.00\,\%  & 0.01\,\%  & 0.18\,\%  & 0.01\,\%  & 0.01\,\%  \\
  & Sample 2 & 11.56\,\% & 6.40\,\%  & 4.53\,\%  & 13.20\,\% & 7.44\,\%  & 4.69\,\%  \\
  & Sample 3 & 14.03\,\% & 6.84\,\%  & 6.71\,\%  & 14.55\,\% & 11.37\,\% & 8.58\,\%  \\
\multirow{3}{*}{5\,\% noise}
  & Sample 1 & 0.36\,\%  & 0.22\,\%  & 0.003\,\% & 0.11\,\%  & 0.15\,\%  & 0.003\,\% \\
  & Sample 2 & 15.29\,\% & 11.24\,\% & 8.12\,\%  & 24.44\,\% & 11.96\,\% & 8.61\,\%  \\
  & Sample 3 & 25.00\,\% & 9.20\,\%  & 8.14\,\%  & 19.03\,\% & 13.61\,\% & 10.58\,\% \\
\bottomrule
\end{tabular}
\end{center}

\vspace{1em}

% ---------- Prose: closing of 1% noise paragraph ----------
\noindent
with a 1\,\% noise level. Furthermore, we provided the values of
the regularization factors for each case depicted in Figs.~8--16
in Table~2.

\vspace{0.8em}

% ---------- Case 2 subsection heading (italic, matching paper style) ----------
\noindent\textit{Case 2: A layered ring structure}

\vspace{0.5em}

\noindent
In the second numerical example, we focused on a layered ring
structure, as illustrated in Fig.~17. To discretize the problem
domain, we employed 7610 triangular elements. The deformation of
the ring structure was induced by applying inner pressures ranging
from 10\,mmHg to 50\,mmHg, and we were only able to ``measure'' the
displacement on the inner wall. The shear modulus values of the
inner and outer layers were 0.129\,MPa and 0.386\,MPa [37],
respectively. Our approach for the inverse problem relied on
utilizing the measured displacement on the inner wall to determine
the shear modulus distribution. Initially, we provided an initial
guess of the shear modulus distribution, as shown in Fig.~18(a),
where four layers were predefined. The shear modulus values, from
the outermost layer to the innermost layer, ranged from 0.01\,MPa
to 0.04\,MPa. During the inverse problem-solving process, we
optimized both the shape of the interface between neighboring
layers and the shear modulus value of each layer. In the initial
assumption, the ring structure is considered to consist of 4
layers. For the explicit inverse method, each interface is
represented by 14 control points, leading to a total of 56
optimization variables, which includes 4 shear moduli. On the
other hand, the implicit inverse method involves a much larger
number of optimization variables, specifically 4030 variables in
total. Fig.~18 displays the variation of the reconstructed shear
modulus distribution using the explicit inverse method for the
noise free case. Notably, the reconstruction achieved through the
explicit method closely resembled the ground truth. However, the
same cannot be said for the implicit inverse method, where the
reconstructed shear modulus distribution performed significantly
worse than the explicit inverse method even for the noise free
data (see Fig.~19(d)-(f)).

% [page ends here — continues on page 12]

\end{document}
